\documentclass[times,twocolumn,final]{elsarticle}

\usepackage{cag}
\usepackage{framed,multirow}
\usepackage{amssymb}
\usepackage{amsmath}
\usepackage[switch]{lineno}
\usepackage{graphicx}
\usepackage{booktabs}
\usepackage{placeins}
\setlength{\emergencystretch}{2em}
\setcounter{totalnumber}{6}
\setcounter{topnumber}{4}
\setcounter{bottomnumber}{4}
\renewcommand{\topfraction}{0.92}
\renewcommand{\dbltopfraction}{0.95}
\renewcommand{\textfraction}{0.05}
\renewcommand{\floatpagefraction}{0.85}


\journal{Computers \& Graphics}

\begin{document}

\begin{frontmatter}

\title{Timestep-Conditioned Adapter Scaling for Multi-view Diffusion Texture Generation: A Controlled Empirical Study}

% Double-anonymized review: author names, affiliations, and acknowledgements
% (including funding) are provided only on the separately uploaded title page.

\begin{abstract}
Geometry-conditioned multi-view diffusion provides an inference-time strength control, but its practical effect depends on temporal allocation, the conditioning interface, and the evaluated output. We conduct a controlled empirical study of Timestep-Conditioned Adapter Scaling (TCAS), carrying forward the historical middle-high profile C3 without further schedule selection. On 300 held-out objects with a frozen retrained GeoTex checkpoint, C3 improves foreground PSNR by 0.52 dB over fixed low, while competitive fixed and linear controls have better means on all four image endpoints. Matching applied-scale budgets changes the interpretation: C3 is 0.13 dB below fixed mean and 0.56 dB below late-high. Native MV-Adapter studies on 50 Fresh C inputs and 100 GSO objects detect endpoint-specific stage effects but do not detect the primary middle-window benefit over fixed mean. A ten-object native-UV bake pilot finds a small late-high LPIPS reduction of 0.00265 on unseen views; a separate twenty-object GSO evaluation with revised calibration does not detect this benefit. Generated-view and baked-output method rankings remain positively associated, although the preferred point estimates often differ. Fixed-seed and shared-UV diagnostics distinguish group-average gains, objectwise repeatability and incoming-view consistency from fidelity. These findings quantify application-specific effects and limits of geometry-control allocation. They support neither a universally preferred schedule nor physically faithful materials.
\end{abstract}

\begin{keyword}
\KWD
multi-view diffusion \sep 3D texturing \sep geometry adapter \sep residual scaling \sep timestep conditioning
\end{keyword}

\end{frontmatter}

\raggedbottom
\linenumbers

\section{Introduction}
\label{sec:intro}

Multi-view diffusion generates images that are projected onto a 3D mesh to form a texture. Geometry conditions help align the images with a known surface, while a reference image supplies appearance. The strength of geometric injection is an accessible inference-time control. Choosing it is an application decision: better generated-view agreement is useful only to the extent that the resulting textured asset retains appropriate appearance under the evaluation views.

The difficulty is that strength, temporal allocation and the final evaluation endpoint can change together. A schedule compared with a weaker constant may simply inject a larger scalar budget. A wrapper may cap some requested scales, leaving a different applied profile. A generated-view metric can also improve while color errors remain visible, and baking adds sampling and coverage losses. Establishing the value of a setting therefore requires competitive controls and actual textured-mesh evaluation.

We study these questions through Timestep-Conditioned Adapter Scaling (TCAS), an intervention on an existing GeoTex-Adapter residual interface. C3, the low--high--low profile recorded in the original development study, is retained without another selection search. Scheduling strength across time or layers is established prior work, including geometry-conditioning schedules in Scheduled Style Injection (SSI) \cite{kulkarni2026ssi}. Here the contribution is the quantitative evidence for geometry-guided multi-view texturing under stated interfaces and protocols.

The complete comparisons change the original practical recommendation. On 300 held-out objects, C3 improves over fixed low but loses to competitive fixed and linear controls. At equal applied-scale sums, fixed mean and late-high outperform its middle window. Native MV-Adapter evaluations on fifty Fresh C inputs and one hundred GSO objects do not detect the primary middle-window advantage, although some secondary stage effects are measurable. Actual baking finds a small late-high benefit on ten compatible GeoTex objects and no detected LPIPS benefit on a separate twenty-object GSO cohort. These findings delimit when an allocation result has application value.

The contributions are:

(1) \textbf{A controlled assessment of strength and temporal allocation.} Matched comparisons with competitive fixed and linear controls distinguish the recorded C3 intervention from a timing-only interpretation. Equal per-group applied sums identify stage-allocation total effects within the frozen GeoTex pipeline.

(2) \textbf{An application-level evaluation of those effects.} Actual native-UV meshes, matched-resolution GT self-bakes and eleven unseen cameras quantify small local gains, persistent appearance errors and changes in the preferred operating point. Both the ten-object gain and the twenty-object non-detection are reported. Known-color controls check the incoming-view disagreement diagnostic, which remains distinct from fidelity.

(3) \textbf{An evaluation of interface and protocol boundaries.} Native MV-Adapter matrices and retained MVDiffusion results test related allocation questions on different conditioning operators. They reveal endpoint-dependent effects and limits to transfer; fixed-seed sensitivity separates an aggregate effect from objectwise repeatability. Cohorts are not pooled and identical scale labels are not treated as a common residual dose.

Object-level records and runnable analyses support these findings. TCAS remains the study intervention and C3 a historical operating point; neither the scheduling concept nor the existence of an audit archive is presented as an independently new algorithm.

\section{Related Work}
\label{sec:related}

\subsection{3D Texture Generation and Multi-view Diffusion}
\label{subsec:3dtex}

Single-image 3D generation is often decomposed into geometry generation and texture generation. Single-image-to-3D, multi-view diffusion, and sparse-view reconstruction methods have improved structural stability and view synthesis \cite{liu2023zero1to3,shi2023zero123pp,liu2023syncdreamer,long2023wonder3d,shi2024mvdream,wang2023imagedream,xu2024instantmesh,ke2024era3d}. Texture synthesis remains a distinct challenge for asset realism and usability \cite{perla2026survey}. Many systems generate RGB views with 2D diffusion and project or bake them onto a surface \cite{chen2023text2tex,richardson2023texture,poole2023dreamfusion,chen2023fantasia3d}. Recent work also explores exemplar-based texture synthesis and diffusion-based procedural materials \cite{guzman2026adaptive,lv2025interpretable}.

Paint3D, MVPaint, Hunyuan3D 2.0, and MVPainter use geometry-conditioned multi-view generation in 3D texturing pipelines \cite{zeng2024paint3d,cheng2024mvpaint,lai2025hunyuan3d,shao2025mvpainter}. MV-Adapter, MVDiffusion, and other synchronized multi-view methods address consistency through different adapter and denoising designs \cite{huang2024mvadapter,tang2023mvdiffusion,liu2024synchronized,xu2026wondertex}. These systems establish the graphics context for this work. The present evidence does not compare final baked-texture quality across those systems.

\subsection{Adapter-Augmented Diffusion Models and Geometric Control}
\label{subsec:adapter}

ControlNet and T2I-Adapter provide learned spatial-conditioning branches or feature injections, whose contribution can be weighted by an inference-time condition-strength scale \cite{mou2023t2iadapter,zhang2023controlnet}. They establish conditioning strength as a practical control, but their injected tensors and scaling sites do not share a numerical dose with GeoTex. The studied GeoTex-Adapter path already exposes an inference-time residual scale and evaluates fixed requests. TCAS reuses that control point as a three-stage request over sampling progress, while recording how wrapper-group caps change the applied values and how adapter activations determine the resulting corrections. It adds no learned branch or schedule selector. IP-Adapter introduces image prompts through decoupled attention \cite{ye2023ipadapter}; LoRA and related methods use scaled low-rank updates \cite{hu2022lora,liu2024dora,zhang2023adalora}.

\subsection{Diffusion Guidance Scheduling and Timestep Control}
\label{subsec:guidance}

Diffusion denoising is often discussed in terms of coarse-to-fine evolution, and prior work varies guidance strength over sampling progress to study quality, diversity, or fidelity \cite{ho2020ddpm,song2021ddim,ho2022cfg,castillo2023adaptive}. Limited-interval classifier-free guidance is one example; it changes the combination of conditional and unconditional score predictions \cite{kynkanniemi2024limited}.

SSI is the closest scheduling precedent. It varies StyleID's query-mixing parameter across decoder layers and denoising time, tests schedule direction and shape, and applies scheduling to a separate geometry-conditioned ControlNet branch; its reported ControlNet scheduling gains over fixed conditioning are modest relative to its style-injection results \cite{kulkarni2026ssi}. This prior work already covers the general idea of layer/time scheduling and geometry-conditioning schedules. TCAS acts on the residual correction produced by GeoTex-Adapter from normal, depth, and foreground-mask conditions, with group-specific wrapper caps; it does not modulate StyleID attention mixing or add a ControlNet branch. The task, interface, and measured endpoints differ, so SSI does not imply the Fresh C paired estimates, the requested/applied/residual accounting, or the observed image- and color-agreement failures. These task-specific measurements are the empirical contribution; the scheduling concept itself is not claimed as new.

\subsection{Multi-view Adapter Interfaces}
\label{subsec:mvinterfaces}

MV-Adapter uses a condition guider for camera or geometry maps together with decoupled spatial, multi-view, and optional image cross-attention layers \cite{huang2024mvadapter}. MVDiffusion uses correspondence-aware attention modules to connect pixels across views \cite{tang2023mvdiffusion}. These interfaces differ from the GeoTex-Adapter residual branch studied here. We therefore report each auxiliary intervention within its own tested implementation and cohort, without pooling metrics or interpreting it as an architecture-matched replication of TCAS. The native temporal comparisons are summarized in Tables~\ref{tab:new-mva} and \ref{tab:interface-effects}; historical panels remain in Supplementary Section~S3.

\section{Study Design and Control Interface}
\label{sec:method}

\subsection{Task Definition and Base Pipeline}
\label{subsec:task}

This paper builds on the MVPainter-style geometry-controlled multi-view texture generation framework \cite{shao2025mvpainter} and studies how the inference-time strength of geometric condition injection affects the generated results. Figure~\ref{fig:overview} shows the TCAS control gate on the GeoTex-Adapter residual path: the requested scale is limited by the wrapper's group-specific cap before the residual is added to frozen UNet features. Given a single reference image and the corresponding 3D geometry, the base pipeline generates RGB texture images from six target views and can pass them to downstream projection or texture baking. Fresh C evaluates generated target views; the ten-object GeoTex pilot and twenty-object GSO evaluation separately assess unlit baked meshes. We do not modify the base UNet, VAE, or pretrained adapter weights.

\begin{figure*}[t]
\centering
\includegraphics[width=\textwidth]{figures__fig1_tcas_control.pdf}
\caption{TCAS on the existing GeoTex residual path. Reference-image and five-channel geometry conditioning enter separate paths. The requested C3 profile is $(1.25,2.50,1.25)$ over steps 0--16, 17--32 and 33--49; native group caps act before residual injection into the frozen UNet. Fresh C shallow requests are capped at 0.8. This editable vector schematic abstracts multiple ResNet sites and is not an experimental result or end-to-end application demonstration.}
\label{fig:overview}
\end{figure*}

During multi-view diffusion, the 3D geometry is rendered into normal maps, depth maps, and foreground masks, denoted as:
\begin{equation}
\mathbf{G} = \{G_\text{normal}, G_\text{depth}, G_\text{mask}\}.
\end{equation}

The geometric adapter maps these conditions to a residual correction term and injects it into the intermediate hidden states of the UNet. For inference step $j$ with normalized execution progress $p_j$, the adapter injection can be written as:
\begin{equation}
\mathbf{h}'_{j,\ell} = \mathbf{h}_{j,\ell} + s_{\mathrm{app}}(p_j,\ell)\, A^\ell_\phi(\mathbf{h}_{j,\ell}, \mathbf{G}),
\label{eq:injection}
\end{equation}
Here $\ell$ indexes an injected layer group. The \emph{requested scale} $s_{\mathrm{req}}(p_j,\ell)$ is the value assigned by the profile; the \emph{applied scale} is the value used by the wrapper,
\begin{equation}
s_{\mathrm{app}}(p_j,\ell)=\min\!\bigl(s_{\mathrm{req}}(p_j,\ell),c_\ell\bigr),
\end{equation}
when the native cap $c_\ell$ is active, and otherwise equals the request. The resulting correction is $\mathbf{r}_{j,\ell}=s_{\mathrm{app}}(p_j,\ell)A^\ell_\phi(\mathbf{h}_{j,\ell},\mathbf{G})$; its \emph{post-scale residual magnitude} is $m_{j,\ell}=\|\mathbf{r}_{j,\ell}\|_2$. Thus neither a requested nor an applied scale label alone determines the correction magnitude.

In the audited Fresh C wrapper, the three cap groups map the GeoTex residual wrappers on the UNet ResNet blocks as follows: deep comprises the UNet middle block and $\texttt{up\_0}$, middle is $\texttt{up\_1}$, and shallow is $\texttt{up\_2}$. These names identify the implementation's grouped injection points; they do not assert universal semantic roles for denoising stages. The corresponding Fresh C caps are 3.0, 3.5, and 0.8. C3's shallow request of 1.25 is therefore applied as 0.8, while its deep and middle requests pass unchanged. GFL and GFH also clip at the shallow cap. These cap values describe the audited Fresh C wrapper; row-level cap semantics for the historical 24-object probe are not fully recoverable. TCAS scales a geometry residual added to feature activations; it does not reweight conditional and unconditional score predictions as classifier-free guidance does.

\subsection{Image-Space Metrics and Texture-Variation Diagnostics}
\label{subsec:metrics}

The endpoints describe different outcomes. \emph{Image-space agreement} refers to target-view PSNR, LPIPS, and SSIM measurements on the named image regions; it is not a measure of final 3D material quality. \emph{Reference-color agreement} refers to color differences such as CIEDE2000 computed against the corresponding reference. \emph{Texture variation} refers to statistics such as RGB standard deviation, Laplacian variance, gradient magnitude, and high-frequency energy; these may rise with detail or artifacts and are not fidelity measures. \emph{Subjective preference} refers only to the submitted forced-choice study. \emph{Baked 3D appearance} refers to rendered mesh outputs after texture projection or baking, while \emph{PBR/material fidelity} additionally concerns appearance under the material and illumination model. These endpoints are not interchangeable. FG-SSIM, Edge-SSIM, and Crop-SSIM quantify structural similarity on the indicated image regions; PSNR measures pixel error; LPIPS compares learned image features \cite{wang2004ssim,zhang2018lpips}.

Laplacian variance, RGB standard deviation, and gradient magnitude can describe local high-frequency or color variation. For example, a foreground Laplacian statistic may be written as
\begin{equation}
M_{\mathrm{var}}=\operatorname{Var}\!\left(\operatorname{Lap}(I_{\mathrm{out}}\odot G_{\mathrm{mask}})\right),
\end{equation}
where $G_{\mathrm{mask}}$ is the foreground mask. A larger value is not necessarily better: detail, noise, repeated patterns, or artifacts can all increase local variation. We therefore do not interpret these proxies as texture fidelity or use them alone to claim that a schedule better preserves reference appearance.

\subsection{TCAS Control Semantics and Comparison Design}
\label{subsec:tcas}

TCAS replaces a globally fixed requested adapter scale with a piecewise profile over normalized sampling progress. The early/middle/late partition is an operational description, not a universal claim about the function of every denoising stage. C3's low--high--low profile is the specific development choice evaluated in this revision; it is not presented as an independently derived optimum.

To avoid ambiguity caused by different raw timestep directions in different schedulers, $p_j=j/(T-1)$ denotes execution progress for inference step $j\in\{0,\ldots,T-1\}$. Thus $p=0$ is the first update from the initial noisy latent and $p=1$ is the last denoising update before decoding. For $T=50$, TCAS divides execution into steps 0--16, 17--32, and 33--49 (17/16/17 steps):
\begin{equation}
s_{\mathrm{req}}(p) = \begin{cases} s_e, & 0 \leq p < 1/3 \\ s_m, & 1/3 \leq p < 2/3 \\ s_l, & 2/3 \leq p \leq 1 \end{cases}.
\label{eq:tcas}
\end{equation}

We call this variation across the 50 inference steps the \emph{temporal allocation}. A \emph{layer allocation} varies the requested scale across the three wrapper groups at a fixed step. For a schedule with applied values $s_{\mathrm{app}}(j,\ell)$, its per-group applied-scale sum is $B_\ell=\sum_j s_{\mathrm{app}}(j,\ell)$; the unweighted sum $B_{\mathrm{nom}}=\sum_\ell B_\ell$ is a \emph{nominal applied-scale budget}. It does not weight scales by adapter activations. We log post-scale correction summaries, including wrapper-level $L_2$ norm and mean absolute entry. The three wrapper groups operate on feature maps with different spatial dimensions, and the traces do not retain channel counts; raw $L_2$ values are therefore not used for cross-layer dose comparisons. The per-entry mean absolute correction reported in Supplementary Section~S1 is a dimension-normalized descriptive summary, not an equal-dose intervention. Correction magnitudes also depend on the current hidden state and earlier denoising updates, so equal nominal budgets do not imply equal realized corrections.

The selected C3 configuration used in this paper is:
\begin{equation}
(s_e, s_m, s_l) = (1.25, 2.50, 1.25).
\end{equation}

The low--high--low profile is a concrete inference-time intervention motivated by the submitted development analysis. Its stage interpretation is a task-specific hypothesis: the tested early, middle, and late allocations may produce different responses. Fresh C evaluates the recorded profile against matched controls but does not prove a universal function for any stage. TCAS adds no trained parameters and does not modify base-model or adapter weights; it replaces the requested scale during inference.

\paragraph{What C3 changes and what its comparisons identify}
C3 requests $(1.25,2.50,1.25)$ for each of the 50 steps across the three layer groups before native caps. In Fresh C, the applied deep and middle profiles retain those values; the shallow profile is clipped to 0.8 throughout. GFL requests 1.25 at every step and group, so its deep and middle values are 1.25 throughout and its shallow value is also clipped to 0.8. Their per-group applied-scale sums are therefore 82.5 versus 62.5 in deep and middle, and 40 for both in shallow. The C3--GFL contrast changes the middle-third deep/middle scale magnitude as well as its temporal allocation; it is not a timing-only comparison.

To identify a temporal-placement contrast, controls must hold the per-layer applied values and nominal per-group sums fixed while moving the same values between time windows. To identify a scale-strength contrast, controls must hold the layer and temporal allocations fixed while changing amplitude. Both contrasts also require reporting post-scale correction summaries: matching requested values or applied-scale sums does not match the evolving residual corrections. Raw $L_2$ values are dimension-sensitive across wrapper groups and cannot establish cross-layer dose matching. The current 300-object allocation comparison exactly matches the applied sums in each layer group. This controls the scalar budget for the total allocation effect, although it does not identify a residual-dose-independent mechanism. Post-hoc adjustment for state-dependent corrections cannot by itself recover that mechanism. The older strict-276 analysis is reported separately.

\paragraph{Development origin and scope of C3}
The submitted 01549 paper records C3 $=(1.25,2.50,1.25)$ as the selected profile from a 24-object development probe containing 17 uniform, layer-wise, and time-conditioned variants. We retain that development record because it explains the origin of the profile studied here. The complete row-level outputs and selection-time provenance are not available in the current evidence authority, so the historical ranking and labels cannot be independently checked. We therefore make no formal optimality or uniqueness claim based on this probe. Fresh C evaluates the recorded profile without further schedule search.

\subsection{Geometry Adapter and Inference Interface}
\label{subsec:adapterimpl}

GeoTex-Adapter is treated as the fixed geometry-residual backbone, not as a new contribution of this paper. It receives a five-channel input comprising a three-channel normal map, a one-channel depth map, and a one-channel foreground mask. A multi-scale encoder produces residual features that are injected into the frozen UNet. TCAS acts on the requested multiplier along denoising progress; it does not retrain the adapter, the UNet, or the VAE and introduces no learned parameters.

Previously reported Full Adaptive Correction (FAC) training is retained as an auxiliary negative result in Section~\ref{subsec:facreported}. It is distinct from the frozen inference interventions evaluated here.

The recorded analysis and display code uses fixed object-level inputs and the original comparison families. Actual experimental images are preserved without appearance retouching. Software identities and the AI-assistance disclosure are supplied with the reproduction materials. Final displays are checked against source hashes, shared cameras/crops and object-level metrics; visual assessment illustrates appearance limitations and is not an additional quality endpoint.

\section{Experiments}
\label{sec:experiments}

\subsection{Study Cohorts, Sampling and Statistical Units}
\label{subsec:setup}

Table~\ref{tab:study-design} separates the studies and their evidence classes. Fresh C contains 300 UIDs disjoint from the historical schedule-selection/evaluation pool, the recorded 1,118-object training list and the earlier clean-v2 evaluation list. Its primary four-condition comparison uses one frozen controlled-retraining checkpoint, shared inputs and per-object seeds. It does not reproduce the historical adapter's reported $+0.96$ dB estimate. The original checkpoint's exact identity remains unavailable.

\begin{table*}[t]
\centering\scriptsize\setlength{\tabcolsep}{3pt}
\caption{Studies, implemented protocols and evidence scope. Rows are studies rather than disjoint populations: the bake pilots and some interface tests reuse image-study inputs. Object is the statistical unit; views and seeds are nested within object.}
\label{tab:study-design}
\begin{tabular}{p{0.16\textwidth}p{0.23\textwidth}p{0.20\textwidth}p{0.32\textwidth}}
\toprule Study / objects & Interface and sampling & Endpoint & Selection and scope \\
\midrule
Fresh C / 300 & Capped GeoTex residual; 256 pixels; 50 Euler steps; conditional-only prediction; cached image features & Six target-view foreground endpoints & UID-disjoint held-out cohort; original primary C3--GFL frozen; linear and equal-budget comparisons are same-cohort addenda \\
GeoTex bake / 10 & Same controlled generator; six 256-pixel bake cameras; native UVs & Eleven unseen unlit views & Input/compatibility-selected subset; primary Late--Mean specified before new bake outcomes; two later native controls exploratory \\
MV-Adapter / 50 & Native condition-guider features; 512 pixels; 50 steps; CFG 3 & Six native-camera target views & Fresh C inputs not previously viewed through this interface; input-only selection; thirteen legal white targets retained \\
GSO / 100 & Native MV-Adapter; 512 pixels; 50 steps; CFG 3; three seeds & All-six PSNR primary; five nonreference LPIPS secondary & Source-ordered convenience cohort; four inherited profiles frozen after earlier baseline exposure; external sensitivity \\
GSO bake / 20 & Same fixed GSO inputs; native UVs; three seeds & Eleven unseen unlit views & Revised GT-only calibration after fifteen original outcomes; full twenty retained; no GT-error subtraction \\
Reference-protocol / 50 & New frozen GeoTex checkpoint; native reference conditions; 512 pixels; 50 steps; CFG 2 & Five nonreference views & Supporting protocol/checkpoint sensitivity on previously studied inputs; not another independent backbone \\
MVDiffusion / 75 & Partial correspondence-block control; depth concatenation unscaled & Recorded six-view endpoints & Retrospective third-interface panel; narrow native-mirror identity check; not exact GeoTex operator transfer \\
\bottomrule
\end{tabular}
\end{table*}

Within each comparison, object inputs, targets, initialization and endpoint code are shared. Fresh C uses seed $42+\texttt{object\_idx}$, held constant across conditions for that object. Its capped residual groups have deep/middle/shallow limits of 3.0/3.5/0.8. MV-Adapter uses a different feature injection operator and numerical scale range; its native target cameras and actual default-path identity are checked separately. Scale labels are not common physical doses across interfaces.

Fresh C's controlled sampler differs from the complete base pipeline, which uses online image features, classifier-free guidance and 512-pixel views. Constant 1.0 under the GeoTex caps also differs from the unset-scale residual branch. We therefore name these controls by their executed configuration. The source records and sampling/default-path diagnostics are in Supplementary Sections~S1 and S6.

FG-PSNR, LPIPS, SSIM and Edge-SSIM measure agreement on the named image regions. Baked evaluation renders actual exported meshes; unlit base color does not establish PBR or material fidelity. Fresh C image endpoints are evaluated once on each frozen $3\times2$ six-view composite. Native-interface and baked endpoints average their recorded per-view scores within object, with GSO sampling seeds also averaged within object. Paired intervals resample objects and are pointwise 95\% intervals. Preserved primary, secondary and exploratory families use their stated Holm adjustments; the distinct families do not imply one overall experiment-wide error guarantee. Non-detection is not equivalence or non-inferiority. Complete comparison rules are supplied in Supplementary Sections~S2--S5.

\subsection{Development origin and historical scope}
\label{subsec:c3}
The historical study selected C3 from 17 profiles on a 24-object probe and reported a separate fixed-scale sweep. Those records explain the intervention's origin, but their complete row-level provenance cannot be reconstructed. They remain historical diagnostics in Supplementary Section~S7. The current matched evaluation carries C3 forward without another selection search. Historical pooled gains, preference and FAC results remain separate from the new evidence.

\subsection{Competitive Controls and Matched Temporal Allocation}
\label{subsec:largescale}

The internally specified Fresh C primary is C3 minus Fixed Low (GFL) FG-PSNR. The paired gain is $+0.52$ dB (95\% interval $[+0.40,+0.64]$), with 206/300 objects favorable. All conditions use the same frozen retrained checkpoint and object inputs. This gain combines increased strength and temporal allocation: C3's deep/middle applied sums are 82.5 each, versus 62.5 for GFL, while both have shallow sum 40.

\begin{table*}[t]
\centering\scriptsize\setlength{\tabcolsep}{4pt}
\caption{Fresh C competitive-control means on the same 300 objects. Early/Middle/Late-high and Fixed Mean share exactly the same per-group applied sums. The equal-budget study is a post-inspection addendum, not a second holdout. Each budget entry for deep/middle applies separately to both groups. $^\dagger$Generic Linear is a competitive same-cohort addendum, not part of the matched-budget family; its different group budgets are shown explicitly.}
\label{tab:freshc-absolute}\label{tab:matched-budget300}
\begin{tabular}{lrrrrrr}
\toprule Control & Deep/middle sum & Shallow sum & FG-PSNR$\uparrow$ & FG-LPIPS$\downarrow$ & FG-SSIM$\uparrow$ & Edge-SSIM$\uparrow$ \\
\midrule
\input{tables__matched_controls014.tex}
\bottomrule
\end{tabular}
\end{table*}

Competitive controls change the recommendation. Fixed High has better means than C3 on all four endpoints: C3 minus Fixed High is $-0.748$ dB FG-PSNR (pointwise interval $[-0.885,-0.609]$), with FG-SSIM difference $-0.04553$ ($[-0.04937,-0.04167]$). Generic Linear ramps deep/middle requests from 1.25 to 2.50 and shallow requests from 0.50 to 0.75 over the fifty steps. Its applied sums are (93.75,93.75,31.25), so it changes layer allocation and total budget as well as timing. It also has better means on all four; C3 minus Generic Linear is $-0.593$ dB ($[-0.837,-0.351]$). Its positive median of $+0.074$ dB and 155/300 favorable objects show why mean and object-level distributions must both be considered. Complete original paired endpoints and their comparison families are in Supplementary Section~S2.

\paragraph{Equal applied budgets.}
\label{subsec:matched-budget300}
Early-, middle-, and late-high move the same sixteen high steps to updates 0--15, 17--32, and 34--49. Deep/middle low and high are 1.25 and 2.50; shallow stays at its 0.8 cap. Fixed Mean uses 1.65 before the cap. Each allocation therefore has per-group sums $(82.5,82.5,40)$. The completed six-condition matrix contains 1,800 paired rows with common-input and seed checks.

C3 minus Fixed Mean is $-0.133$ dB (pointwise interval $[-0.170,-0.096]$); C3 minus Late-high is $-0.564$ dB ($[-0.634,-0.494]$). C3 also trails both on LPIPS, FG-SSIM and Edge-SSIM. It improves over Early-high by approximately 0.50 dB. Within these six conditions Fixed High wins FG-PSNR on 172 objects, Late-high on 64, and C3 on none; this is an endpoint-specific count, not a universal quality ranking.

A descriptive four-endpoint comparison complements these winner counts. On 270/300 objects, at least one of the other five controls weakly dominates C3 on all four recorded composite endpoints and improves at least one. Fixed High alone does so on 173 objects, but C3 dominates it on 16; their remaining 111 objects have competing endpoints. This post-inspection, point-estimate analysis is not statistical separation on each object or a ranking of unmeasured material quality. Full counts are in Supplementary Section~S2.

The equal-budget comparison identifies an allocation total effect. Evolving residual amplitudes remain trajectory-dependent, preventing a residual-dose-independent mechanism interpretation. Late-high is not universally preferred: Fixed Mean has higher FG-SSIM, and Generic Linear is a strong unmatched-budget practical control. The historical LLH layer/time addendum and complete object distributions remain in Supplementary Section~S2.

\begin{figure*}[!t]
\centering
\includegraphics[width=\textwidth]{figures__fig5_freshc_context_roi.pdf}
\caption{Fresh C view 0: (a) Octa\`edre, a fixed-scale endpoint trade-off; (b) Cube 01, a Generic Linear counterexample. Each column shows the complete object in a common frame and the same pixel ROI enlarged below. Object-level C3 minus Generic Linear in (b) is $-3.491$~dB FG-PSNR and $+0.08172$ FG-LPIPS, favoring Linear on both. Added patterns visibly differ from the low-variation target. These retained outcome-selected examples illustrate behavior, not prevalence or material fidelity. Composite object metrics use all six views, not only the displayed crop. Source hashes, camera/seed identities and display frames are in the final visual manifest; credits are in Supplementary Section~S9.}
\label{fig:c3compare}
\end{figure*}


\subsection{Appearance and Color Failure Analysis}
\label{subsec:appearance-color}

The selected color diagnosis contains 26 failure-enriched objects: six earlier figure cases and 20 Fresh C panels. It is not a random sample, a new holdout, or a population estimate. No Adapter, GFL, and LLH were compared on the same inputs, checkpoint, and runner with six target views per object (78 object-condition records). The following paired object-bootstrap intervals are descriptive (20,000 resamples; seed 61008), and no confirmatory multiplicity family was defined.

\begin{figure*}[t]
\centering
\includegraphics[width=\textwidth]{figures__fig6_color_diagnostic_freshc_context_roi.pdf}
\caption{Fresh C view 0, E5-5-2: target, Fixed Low, C3 and layer LLH share a complete-object frame and an identical pixel crop. The generated brown pattern does not reproduce the gray target; changing its variation does not establish color fidelity. This previously viewed example came from an existing C3--GFL LPIPS stratum and was not selected for color. It does not reproduce the legacy purple-shift cases. The malformed disabled-adapter panel is omitted from this display; its numerical results remain in the full analysis. Source/display identities and attribution are supplied with the final visual manifest and Supplementary Section~S9.}
\label{fig:color-failure}
\end{figure*}

\begin{table}[t]
\centering
\scriptsize
\setlength{\tabcolsep}{3pt}
\caption{Selected paired results from the failure-enriched 26-object color diagnosis. Negative CIEDE2000 differences indicate lower reference-color error for the first condition. The set is descriptive and does not estimate population prevalence.}
\label{tab:color-summary}
\begin{tabular}{p{2.3cm}p{3.3cm}p{1.3cm}}
\toprule
Contrast & Mean CIEDE2000 difference [95\% CI] & Objects \\
\midrule
GFL -- No Adapter & $-19.654$ [$-26.370,-13.248$] & 24/26 lower \\
LLH -- GFL & $+3.583$ [$+0.462,+6.712$] & -- \\
\bottomrule
\end{tabular}
\end{table}

GFL lowered reference-relative CIEDE2000 in 24/26 objects, while visible pink or purple shifts remained in two legacy Fig.~4 cases under No Adapter, GFL, and LLH. Their mapping to the original submitted panels cannot be recovered; the traceable Fresh C example in Figure~\ref{fig:color-failure} is a different object and is not presented as a reproduction of those cases. The decoded RGB output showed the shifts before UV projection or baking, and a late denoising snapshot (update 40 of 50) showed one of them. In that case, a ground-truth VAE encode--decode isolate had CIEDE2000 4.086 and $\Delta a^*\approx-1.35$, compared with 22.032 and $+1.977$ for GFL. This weakens ordinary VAE round-trip error as the sole explanation in that case; it does not identify the complete cause or rule out other model and conditioning effects.

LLH--GFL increased mean CIEDE2000 while reducing foreground gradient magnitude, high-frequency energy and color entropy. Lower variation therefore did not imply closer reference appearance. The denoising/runtime isolates and three failed correction candidates remain in Supplementary Section~S6; they localize observations but do not validate a cause or repair.

Historical twenty-object results, the earlier five-object bake, invalidated camera/mask attempts and the corrected REPEAT sampler are reported separately in Supplementary Section~S5. They are not pooled with the new ten-object cohort. The actual unset-scale GeoTex residual branch differs from Constant 1.0 (capped), so the latter is not labeled an unmodified pipeline.

\paragraph{Full-object evidence and its selection.}
The strong-fixed and Generic Linear examples in Figure~\ref{fig:c3compare} and the color failure in Figure~\ref{fig:color-failure} use common framing and exact shared pixel crops. Favorable/adverse C3--GFL examples are retained in Supplementary Section~S9. Outcome-selected displays illustrate behavior, not prevalence. The unrecoverable historical texture audit's 92\% prevalence claim and its texture-loss classifications are withdrawn; variation statistics alone would not validate fidelity even if their lineage were complete.


\subsection{Actual Baked Appearance: Local Benefit and External Limits}
\label{subsec:new-bake}
We froze ten previously unbaked Fresh C objects in UID order using input availability, native UV/material compatibility, triangle-count and camera-mask alignment criteria. Of the 300 screened inputs, only ten passed all criteria before the cutoff; the planned maximum was twenty. The ten are not new to the GeoTex image domain, but their baked outcomes had not been inspected during selection. This is a compatibility-selected pilot, not a representative twenty-object confirmation. Five earlier bake objects are reported separately.

Seven conditions share each object's source mesh, reference, seed, six projection cameras, masks, 256-pixel source resolution, atlas resolution, hole filling and eleven unseen evaluation cameras. Original geometry, UVs and legal base-color material semantics are retained and checked. Unseen cameras do not enter the bake. GT-256 uses the same source resolution, masks and bake parameters and has exactly the same raw atlas coverage as every generated condition. The separately tested CPU REPEAT sampler passes boundary and integer-period tests; it is not the sampler used by this GPU bake. Pole-view RGB orientation is corrected without rotating raw-camera masks. These calibrations control projection artifacts; evaluation is unlit base color, not PBR fidelity. Supplementary Section~S5 records the checks and object-level data.

\begin{table}[ht]
\centering\scriptsize
\caption{New ten-object unlit unseen-view bake pilot. Eleven views are averaged per object. Higher PSNR/SSIM and lower LPIPS/CIEDE2000 are favorable. Constant 1.0 is an explicitly capped request, not the unset-scale pipeline. All seven original generated conditions and all ten eligible objects are retained; the two subsequently added controls marked $\dagger$ are exploratory; GT-256 is a separate matched reference bake, not a theoretical bound.}
\label{tab:new-bake}
\begin{tabular}{lrrrr}
\toprule Condition & FG-PSNR & FG-LPIPS & FG-SSIM & CIEDE2000 \\
\midrule
GT-256 reference bake & 22.38 & 0.0590 & 0.7683 & 4.44 \\
\midrule
Native residual$^\dagger$ & 11.13 & 0.1372 & 0.3376 & 26.47 \\
No Adapter$^\dagger$ & 5.92 & 0.1600 & 0.3829 & 40.92 \\
Constant 1.0 (capped) & 11.93 & 0.1354 & 0.3504 & 24.18 \\
Fixed Low & 12.48 & 0.1317 & 0.3835 & 22.71 \\
Fixed High & 12.29 & 0.1305 & 0.4587 & 23.27 \\
Fixed Mean & 12.62 & 0.1303 & 0.4193 & 22.32 \\
C3 / Middle-high & 12.48 & 0.1302 & 0.4059 & 22.16 \\
Late-high & 12.92 & 0.1277 & 0.4259 & 21.65 \\
Generic Linear & 11.60 & 0.1276 & 0.4458 & 24.71 \\
\bottomrule\end{tabular}
\end{table}

The internally predeclared primary contrast was Late-high minus Fixed Mean on unseen-view FG-LPIPS, proposed from the earlier pilot and tested on the ten newly selected bake objects. Its mean is $-0.00265$ (95\% object-bootstrap CI $[-0.00414,-0.00137]$; exact paired sign-flip $p=0.00195$), with all ten objects favorable. This is approximately a 2.0\% reduction relative to Fixed Mean's mean LPIPS, and identifies a small allocation benefit within this sampling protocol. Generic Linear has a slightly lower mean LPIPS but lower PSNR; Late-high does not dominate every endpoint or control. The matched GT-256 reference bake has mean LPIPS 0.0590, versus 0.1303 for Fixed Mean and 0.1277 for Late-high. A post hoc description therefore puts the absolute gain in context: it closes only 3.7\% of the mean Fixed Mean--GT reference gap (pointwise object-bootstrap CI [2.3\%,5.3\%]). This ratio of paired means is not a theoretical ceiling, a new primary test or evidence of acceptable material fidelity. Leaving out each object in turn keeps the mean Late--Mean difference between $-0.00294$ and $-0.00208$, but all substantial reference-appearance discrepancies remain. Supplementary Section~S5 supplies the paired values and calculation.

C3--Fixed High has mean PSNR $+0.19$~dB (95\% CI $[-0.65,+0.87]$) and LPIPS $-0.00036$ ($[-0.00396,+0.00336]$); neither separates under the four-comparison secondary Holm family (both adjusted $p=1$). Its SSIM is lower by 0.05274 ($[-0.07459,-0.03170]$; Holm $p=0.00781$). Thus the old five-object C3 PSNR advantage does not repeat stably on this new pilot. We retain C3 as the historical candidate and do not rename Late-high as C3.

Across the seven conditions, mean within-object Spearman correlations between generated-view and baked unseen-view rankings are 0.846 for PSNR, 0.804 for LPIPS, 0.746 for SSIM and 0.907 for Edge-SSIM. Broad rankings are positively associated, while the winning condition changes on 7/10 objects for LPIPS (1/10 for PSNR). These are different statements: we do not claim a general reversal of rankings. Full-object Figure~\ref{fig:baked-context} also shows persistent appearance errors, including the Tire's brown output against a gray reference. A lower LPIPS does not establish faithful material reproduction.

\begin{figure*}[t]
\centering
\includegraphics[width=\textwidth]{figures__fig8_baked_unseen_context.pdf}
\caption{Unseen camera 004: (a) LG Remote and (b) Tire. Source and matched GT-256 bake separate projection loss from generated appearance; Fixed Mean and Late-high show the specified allocation comparison. Columns share a full-object frame determined from source, GT and every original condition, without color or detail editing. These are the first two input-order objects with verified publication licenses, not quality-selected examples. Persistent brown-pattern and material errors remain visible. Complete seven-condition displays for all seven rights-cleared objects are in Supplementary Section~S5; all ten objects enter the numerical analysis. Credits: mohd.akmal.mahyudin and Koto3D Stephane Vogley, CC Attribution (full source links in Supplementary Section~S5).}
\label{fig:baked-context}
\end{figure*}


\paragraph{Separate twenty-object GSO evaluation.}
The GSO/MV-Adapter study completes twelve conditions by three seeds on twenty fixed inputs, plus twenty matched GT self-bakes: 740 unique baked meshes and 8,140 unseen-view records. The objects, source meshes, projection cameras, texture resolution and eleven unseen cameras are shared across controls; unseen cameras do not contribute to baking.

The original attempt stopped after fifteen objects because a discontinuous UV-ramp diagnostic failed on a Hair Dryer. Subsequent GT-only tests identify finite-resolution mixing across UV jumps and validate continuous nonlinear surface-color fields, an independent multi-island analytic fixture, and a deliberate image-flip negative control. The resulting calibration addendum is post hoc. It keeps projection, masks, cameras, material semantics, resolutions and the fixed twenty objects unchanged; all 375 original GLB hashes and metric records reproduce exactly. The original failure and remaining UV/coverage limitations are retained in Supplementary Section~S4.

Late-high minus Fixed Mean unseen-view LPIPS is $+0.000084$ (pointwise interval $[-0.000103,+0.000286]$; original comparison-family Holm $p=0.788$). This study does not detect the ten-object pilot's benefit. It changes both input domain and conditioning/projection protocol, so it is neither a direct same-protocol replication nor evidence assigning the difference to backbone alone. Table~\ref{tab:bake-cohorts} presents both application comparisons.

\begin{table*}[t]\centering\scriptsize\setlength{\tabcolsep}{4pt}
\caption{Late-high minus Fixed Mean on unseen-view foreground LPIPS. Lower is favorable. All views/seeds are nested within object. The ten-object specified primary uses its original exact sign-flip test; GSO uses the original Holm16 comparison family. Pointwise intervals are not simultaneous. The two cohorts and protocols are not pooled.}
\label{tab:bake-cohorts}
\begin{tabular}{lrrrrl}
\toprule Study & Objects & Fixed Mean & Late-high & Paired difference [95\% CI] & Test \\
\midrule
GeoTex compatibility pilot & 10 & 0.1303 & 0.1277 & $-0.002652\ [-0.004144,-0.001367]$ & exact $p=0.00195$ \\
GSO revised-calibration study & 20 & 0.10089 & 0.10098 & $+0.000084\ [-0.000103,+0.000286]$ & Holm $p=0.788$ \\
\bottomrule\end{tabular}\end{table*}

Across all twelve GSO conditions, mean within-object image/bake LPIPS-rank correlation is 0.601 (object interval $[0.440,0.733]$). A post hoc sensitivity restricted to the seven original scale/time controls gives 0.396 ($[0.175,0.582]$), with the minimum-mean LPIPS condition changing on 11/20 objects. These are positively associated rankings with changing point-estimate winners, not general rank reversal or statistically resolved winner changes. The GT self-bake has mean LPIPS 0.0602 and PSNR 19.36 dB; average complete-bilinear-footprint coverage is about 78.7\%. All foreground remains scored and GT errors are not subtracted. These losses contextualize the small method effects without defining a theoretical quality bound.

\paragraph{Incoming-view consistency.}
The ten-object bake also records the mean linear-RGB variance across cameras at UV texels receiving at least two visible source views. Replaying the frozen projection code reproduces all 80 object--condition values exactly. On the first two input-order objects, a shared constant-black control gives zero variance; known camera-specific gray levels match an independent analytic variance within $1.1\times10^{-8}$. A wrong constant appearance can therefore be fully consistent by this measure. Late-high minus Fixed Mean is $-0.002448$ (descriptive object interval $[-0.003473,-0.001604]$; additional post-inspection Holm21 $p=0.041$). Fixed High has still lower mean disagreement than Late-high, but a higher baked LPIPS. Resampling, natural texture and overlapping UVs contribute to this diagnostic: it assesses incoming-view disagreement and does not validate seams, a faithful texture or PBR. Formula, all controls and contrasts are in Supplementary Section~S5.

\paragraph{Native controls and protocol sensitivity.}
Subsequently added Native residual and No Adapter outputs on the same ten GeoTex objects are exploratory. They retain matched geometry and GT coverage, and the repeated Constant 1.0 (capped) outputs are bitwise identical. Their apparent PSNR/LPIPS gains for the schedules do not separate under the fifteen-test exploratory family. The complete data are in Supplementary Section~S5.

A supporting fifty-object comparison uses a newly trained frozen GeoTex checkpoint with 512-pixel supervision, actual positive/negative reference conditions and CFG 2. Late-high minus Fixed Mean has nonreference PSNR $-1.184$ dB (interval $[-1.587,-0.849]$) and LPIPS $+0.00266$ ($[-0.00003,+0.00640]$; adjusted $p=0.080$). C3 has no corrected advantage over Fixed Mean or Fixed High. This previously studied cohort and changed checkpoint/protocol supply a sensitivity boundary, not an independent backbone confirmation. Training and full-v29 strong-baseline results remain in Supplementary Section~S8.

\subsection{Transfer to other conditioning interfaces}
\label{subsec:cross-interface}
MV-Adapter and MVDiffusion supply different conditioning operators and scale semantics from GeoTex. We evaluate each within its native target/camera contract and do not pool their scores or interpret identical scale labels as common residual doses.

The historical 76-object MV-Adapter layer/time panel has mixed endpoints and unadjusted temporal intervals. In the 75-object MVDiffusion panel, the layer-wise L--LLH condition is worse than global fixed low in PSNR, structural agreement and reference-color error; LPIPS does not separate. The intervention scales correspondence blocks rather than the same GeoTex residual operator. These counter-results, including all contrasts and identity-check scope, remain in Supplementary Section~S3. They limit transfer claims; low absolute quality is not grounds for excluding them.

\paragraph{New native MV-Adapter temporal comparison.}
\label{subsec:new-mva}
We froze fifty Fresh C inputs not previously inspected through MV-Adapter, excluding all recorded earlier native and profile-panel objects. The input-only ordered selection used availability and native material/geometry compatibility. The native mesh normalization, six orthographic cameras and nvdiffrast rasterizer also generate evaluation targets; an independent source-UV sampling check validates this target path. Three textured new objects pass official/deployed and constant-default-path identity checks: 18 RGB images, conditioning features and injected residuals are bitwise identical. The old seven-object quality archive fails camera calibration and is superseded, not pooled.

The seven conditions share object inputs, seed, geometry and evaluation code. With 50 steps, low/high are 0.75/1.0; each stage-high schedule has 16 high steps, Fixed Mean is 0.83, and Generic Linear is 0.75--0.91. These five mean-matched controls each sum to 41.5 per scaled feature, within this interface. The fifty-object matrix is complete (350 rows), including thirteen texture-absent sources with legal constant-white base-color targets and a documented loader compatibility path. The other 37 textured objects are a descriptive subset, not a replacement primary cohort. Targets are unlit native-camera base color, not Fresh C's differently normalized reference renders or certified physical albedo.

\begin{table}[ht]\centering\scriptsize
\caption{Native-camera MV-Adapter comparison, fifty input-selected objects. Views are averaged within object. Scale values and endpoints are specific to this interface, and are not numerical doses shared with GeoTex.}\label{tab:new-mva}
\begin{tabular}{lrrrr}
\toprule Condition & FG-PSNR & FG-LPIPS & FG-SSIM & Edge-SSIM \\
\midrule
Native / Fixed High & 12.60 & 0.1280 & 0.5041 & 0.3526 \\
Fixed Low & 12.73 & 0.1270 & 0.5066 & 0.3475 \\
Fixed Mean & 12.68 & 0.1279 & 0.5068 & 0.3509 \\
Early-high & 12.60 & 0.1296 & 0.5041 & 0.3503 \\
Middle-high & 12.71 & 0.1262 & 0.5061 & 0.3470 \\
Late-high & 12.74 & 0.1260 & 0.5075 & 0.3514 \\
Generic Linear & 12.72 & 0.1263 & 0.5072 & 0.3501 \\
\bottomrule\end{tabular}\end{table}

The internally predeclared primary, Middle-high minus Fixed Mean FG-PSNR, is $+0.029$~dB (95\% object-bootstrap CI $[-0.019,+0.081]$; paired Monte Carlo sign-flip $p=0.281$; 29/50 favorable). This does not support a stable middle-window gain or establish equivalence. The remaining 104 pair--metric tests use one Holm family. Middle-high has lower Edge-SSIM than Fixed Mean (adjusted $p=0.00104$); some other placement contrasts separate, including Late-high versus Early-high PSNR. Such endpoint-dependent differences do not demonstrate transferable C3 optimality. Complete contrasts, correction rules and the textured subset are in Supplementary Section~S3. This is a native-interface validation of the scheduling question, not architecture-matched replication of GeoTex's capped residual intervention.


\begin{figure*}[t]
\centering
\includegraphics[width=\textwidth]{figures__fig9_mvadapter_native014.pdf}
\caption{Native MV-Adapter view 0 for the same first two input-order rights-cleared objects: (a) LG Remote, (b) Tire. Target, Fixed Mean, Middle-high and Late-high use the native camera, object seed and a common complete-object frame. Original backgrounds and research pixels are retained. The remote preserves recognizable details while the tire retains a marked appearance mismatch; the neighboring scheduling outputs are visually similar. These illustrative cases do not establish cross-interface superiority; the fifty-object primary interval includes zero. All seven controls are displayed in Supplementary Section~S3. Credits: mohd.akmal.mahyudin and Koto3D Stephane Vogley, CC Attribution (Supplementary Sections~S5 and S9).}
\label{fig:native50-visual}
\end{figure*}

\paragraph{External-input native schedule comparison.}
The source-ordered GSO100 convenience cohort uses the same native MV-Adapter control path and three shared seeds. Native/Fixed High, Fixed Low, Fixed Mean, Early-high, Middle-high, Late-high and Generic Linear include the four inherited profiles frozen before their outputs, after earlier baseline exposure. The full dataset contains 3,600 generated cells and 21,600 view records, including three archived projection corrections and two separately frozen APG baselines \cite{sad2024apg}. The scheduling extension is external sensitivity, not pristine independent confirmation or parameter search.

The specified all-six-view Middle-high minus Fixed Mean primary is $-0.014$ dB (interval $[-0.031,+0.004]$; paired $p=0.122$). A secondary Late-high--Fixed Mean contrast is $+0.045$ dB ($[+0.024,+0.065]$; Holm83 $p=0.00335$), with SSIM $+0.00134$ (adjusted $p=0.000415$). Its five-nonreference-view LPIPS difference is $-0.000086$ ($[-0.000551,+0.000413]$; separate Holm21 $p=1$). Thus a small detected image effect does not establish a perceptual or baked-appearance benefit. Table~\ref{tab:interface-effects} and Figure~\ref{fig:allocation-effects} put these named comparisons in context.

\begin{table*}[t]\centering\scriptsize\setlength{\tabcolsep}{3pt}
\caption{Selected within-study paired effects. P: the specified primary, unadjusted within its study; S: secondary with its original family adjustment (GSO all-six Holm83, nonreference LPIPS Holm21, reference-protocol secondary Holm18). Intervals are pointwise. Native50 and GSO100 use different inputs/seed policies; the reference-protocol study also changes weights and sampling and is not a separate backbone.}
\label{tab:interface-effects}
\begin{tabular}{lllrrcr}
\toprule Study & Contrast & Endpoint & Mean & 95\% CI & Class & $p$ \\
\midrule
\input{tables__interface_effects014.tex}
\bottomrule\end{tabular}\end{table*}

\begin{figure*}[t]\centering
\includegraphics[width=\textwidth]{figures__fig_allocation_effects014.pdf}
\caption{Paired allocation effects from completed object-level records. Error bars are pointwise 95\% object intervals. Panel (a) uses Fresh C composite PSNR and native-interface mean-view PSNR; the reference-protocol row excludes its reference view; panel (b) uses eleven unseen baked views. The GSO Late--Mean PSNR row is secondary and is explicitly distinguished from the specified middle-window questions. Different studies are shown separately, without pooling or cross-pipeline absolute-score ranking. The small GeoTex bake gain and GSO non-detection both remain visible.}
\label{fig:allocation-effects}
\end{figure*}

\paragraph{Fixed-seed sensitivity.}
A post-inspection analysis retains all three recorded sampling seeds. GSO100 Late--Mean PSNR averages are positive at each seed ($+0.057,+0.038,+0.039$ dB), but only 27/100 objects improve at all three and 68/100 have mixed or tied directions. Its nonreference LPIPS averages change direction ($-0.000389,-0.000269,+0.000398$). For the GSO20 bake, only 1/20 objects has lower Late--Mean LPIPS at all three seeds; 14/20 have mixed or tied directions. Sign agreement has no practical-effect threshold. These conditional, object-level summaries do not estimate training randomness or revise the original primary tests; full seed intervals and deltas are in Supplementary Section~S4.

All fifteen specified GSO projection-minus-baseline nonreference-LPIPS comparisons favor the scale/APG baselines after their original separate Holm9/Holm6 corrections. The three projection candidates also lose to native on baked LPIPS in the retained twenty-object family. These valid negative auxiliary results are reported completely in Supplementary Sections~S4 and S8; they are not proposed as additional methods.

\subsection{Historical Preference and Training Evidence}
\label{subsec:preference}
The original blinded 3AFC study on thirty objects and twenty-four participants reports C3's overall-quality share as 58.1\% (two-level participant/object-bootstrap interval $[50.1,65.8]\%$). This supports overall preference for the original stimuli, not texture fidelity, the retrained checkpoint or unseen baked appearance. Counts and endpoint-specific results are retained in Supplementary Section~S7.

\paragraph{FAC re-examination.}
\label{subsec:facreported}
Previously reported FAC configurations all reduced FG-PSNR relative to TCAS in their stated controlled re-examination. The original training pool was reported as 1,706 objects, while a recovered configuration refers to 1,200; seed, checkpoint and sampled-run records are incomplete. The results and reproducible inference/statistical components are separated from this remaining training-replay gap in Supplementary Section~S7. Recovered FAC source is now supplied with standalone controller checks. An earlier entry requires an explicit flag to freeze the base adapter, whereas the later entry freezes it by construction; the executed historical source/CLI cannot be authenticated from their names. Later new geometry training does not restore historical FAC identity.


\FloatBarrier
\section{Discussion and Limitations}
\label{sec:limitations}

\paragraph{What the controls establish.}
The Fresh C300 gain over fixed low supports the recorded intervention under its frozen protocol. It changes both budget and timing. The equal-applied-budget matrix then measures a timing-allocation total effect and changes the preferred image-space operating point. It does not identify a state-independent residual dose or an optimal middle window. Competitive fixed and linear controls are necessary before assigning practical value to a selected profile.

\paragraph{Application value and operating points.}
In the GeoTex image protocol, Fixed High and Generic Linear have stronger endpoint means than C3; equal-budget late-high trades higher PSNR and lower LPIPS against lower FG-SSIM than fixed mean. The ten-object bake's approximately 2\% LPIPS reduction closes only 3.7\% of the gap to the matched GT reference bake, while visible appearance errors remain. The separate GSO20 study does not detect that allocation benefit. Its post hoc calibration, GT self-bake loss and incomplete coverage are disclosed; a small image effect alone does not justify replacing an existing deployment setting.

The evidence suggests a practical evaluation sequence rather than a new selection algorithm: compare native/default and competitive fixed settings inside the same deployed protocol; use matched applied budgets to assess timing; check reference appearance and the actual unseen-view bake; and repeat the named comparison across fixed seeds before claiming repeatable objectwise improvement. The tested GeoTex image setting supports stronger fixed controls than historical C3. The native-interface and GSO bake results provide no general reason to replace a deployment's fixed setting with the middle window. These bounded conclusions do not depend on a universally winning schedule.

Generated-view metrics remain informative: overall baked/image method rankings are positively associated. However, preferred point estimates often change, and the association is weaker when the GSO analysis is restricted to closely competing scale/time settings. Winner changes do not prove significant separation between methods or a general reversal. Actual baked evaluation and full-object reference comparison remain necessary.

\paragraph{Interface and protocol boundaries.}
Native MV-Adapter50 and GSO100 do not detect the primary middle-window advantage. The secondary GSO late-window PSNR change is statistically detectable but small, with no corresponding nonreference LPIPS benefit. The supporting reference-conditioned checkpoint shows a different Late--Mean direction. Historical MVDiffusion supplies a partial-operator boundary. Because weights, inputs, conditioning operators and sampling protocols differ, these studies cannot isolate one backbone or domain as the cause of the varying effects.

\paragraph{Quality, selection and replication limits.}
The bake studies favor compatible native-UV assets, and GSO uses a source-ordered convenience sample. They do not establish representative asset quality, physical albedo, PBR fidelity or a validated seam score. Views and seeds are nested within object, and seed repetitions do not estimate training-checkpoint variation. The new seed, dominance and UV-overlap summaries were proposed after observing the original outcomes and remain descriptive; their diagnostic checks do not change that chronology. Separate test families and pointwise intervals retain their stated roles; non-detection is not equivalence.

Fresh C uses retrained weights and a controlled sampler rather than the authenticated historical main adapter or complete default pipeline. The old $+0.96$ dB contrast is not reproduced, and historical FAC training records remain incomplete. The original preference study concerns its own stimuli. These limitations constrain deployment and replication claims, while the matched within-protocol comparisons remain valid.

\section{Conclusion}
\label{sec:conclusion}

Geometry-control allocation changes generated-view agreement, but its effects depend on the tested interface, protocol and endpoint. The matched 300-object study shows that C3's gain over fixed low combines increased strength with timing; equal-budget fixed mean and late-high outperform the historical middle window. Native MV-Adapter studies do not detect the primary middle-window advantage, while some secondary image effects are small but measurable. Actual baking yields a limited late-high benefit on ten compatible GeoTex objects and no detected LPIPS benefit on a separate twenty-object GSO cohort. Method rankings are broadly associated across image and baked evaluation, yet their preferred operating points often differ. Competitive controls and application-level evaluation are therefore necessary when choosing a setting within a deployed pipeline. The results do not identify a universally preferred schedule or establish faithful material reproduction.

\section*{Data availability}

Object-level revision metrics, frozen protocols/configuration records, verified inference and analysis scripts, and the LaTeX source required to rebuild the manuscript and Supplement are included in the author-review package. Recovered FAC source and its source/entry audit are supplied separately; standalone controller tests do not constitute historical full-model replay. The standalone reproduction archive supports statistic recomputation without model weights; fresh inference additionally requires the recorded weights, runtime and licensed source assets. The complete raw generated-image archives, trained checkpoints, and participant-level responses are not included. Source assets must be acquired under their recorded licenses; the review evidence includes attributed displays and a small licensed GSO example export. Additional analysis records may be requested from the corresponding author, subject to licensing, privacy, and archive constraints.

\section*{Declaration of generative AI and AI-assisted technologies in manuscript preparation}
OpenAI Codex assisted with revision drafting, implementation and reproducible statistical and figure code. The final interface diagram uses reproducible vector elements; research displays come from the recorded runs without image-generation retouching. A vision-capable AI assistant inspected final figures and rendered pages for legibility and consistency with the named comparisons; this check does not certify scientific validity or replace author review. Numerical tables are computed from delivered object-level records. Final scientific review, accuracy and responsibility remain with the authors.

\section*{Declaration of competing interests}

The authors declare that they have no known competing financial interests or personal relationships that could have appeared to influence the work reported in this paper.


\begingroup
\linespread{0.94}\selectfont
\begin{thebibliography}{45}
\bibitem{liu2023zero1to3}
Liu R, Wu R, Van Hoorick B, et al. Zero-1-to-3: Zero-shot one image to 3D object. \textit{Proceedings of the IEEE/CVF International Conference on Computer Vision}, 2023.

\bibitem{shi2023zero123pp}
Shi R, Chen H, Zhang Z, et al. Zero123++: A single image to consistent multi-view diffusion base model. arXiv:2310.15110, 2023.

\bibitem{liu2023syncdreamer}
Liu Y, Lin C, Zeng Z, et al. SyncDreamer: Generating multiview-consistent images from a single-view image. \textit{International Conference on Learning Representations}, 2024.

\bibitem{long2023wonder3d}
Long X, Guo Y C, Lin C, et al. Wonder3D: Single image to 3D using cross-domain diffusion. \textit{Proceedings of the IEEE/CVF Conference on Computer Vision and Pattern Recognition}, 2024.

\bibitem{shi2024mvdream}
Shi Y, Wang P, Ye J, et al. MVDream: Multi-view diffusion for 3D generation. \textit{International Conference on Learning Representations}, 2024.

\bibitem{wang2023imagedream}
Wang P, et al. ImageDream: Image-prompt multi-view diffusion for 3D generation. arXiv:2312.02201, 2023.

\bibitem{xu2024instantmesh}
Xu J, et al. InstantMesh: Efficient 3D mesh generation from a single image with sparse-view large reconstruction models. \textit{Advances in Neural Information Processing Systems}, 2024.

\bibitem{ke2024era3d}
Li P, Liu Y, Long X, et al. Era3D: High-resolution multiview diffusion using efficient row-wise attention. \textit{Advances in Neural Information Processing Systems}, 2024.

\bibitem{chen2023text2tex}
Chen D Z, et al. Text2Tex: Text-driven texture synthesis via diffusion models. \textit{Proceedings of the IEEE/CVF International Conference on Computer Vision}, 2023.

\bibitem{richardson2023texture}
Richardson E, Metzer G, Alaluf Y, et al. TEXTure: Text-guided texturing of 3D shapes. \textit{ACM SIGGRAPH 2023 Conference Proceedings}, 2023.

\bibitem{poole2023dreamfusion}
Poole B, Jain A, Barron J T, Mildenhall B. DreamFusion: Text-to-3D using 2D diffusion. \textit{International Conference on Learning Representations}, 2023.

\bibitem{chen2023fantasia3d}
Chen R, Chen Y, Jiao N, Jia K. Fantasia3D: Disentangling geometry and appearance for high-quality text-to-3D content creation. \textit{Proceedings of the IEEE/CVF International Conference on Computer Vision}, 2023.

\bibitem{shao2025mvpainter}
Shao M, Xiong F, Sun Z, Xu M. MVPainter: Accurate and detailed 3D texture generation via multi-view diffusion with geometric control. arXiv:2505.12635, 2025.

\bibitem{kynkanniemi2024limited}
Kynk\"a\"anniemi T, Aittala M, Karras T, et al. Applying guidance in a limited interval improves sample and distribution quality in diffusion models. \textit{Advances in Neural Information Processing Systems}, 2024.

\bibitem{perla2026survey}
Perla S R K, Zhang H, Mahdavi-Amiri A. Advances in neural 3D mesh texturing: A survey. \textit{Computer Graphics Forum}, 2026, e70392.

\bibitem{guzman2026adaptive}
Guzm\'an J E, Mould D, Paquette E. Adaptive multiresolution exemplar-based texture synthesis on animated fluids. \textit{Computers \& Graphics}, 2026, 134: 104490.

\bibitem{lv2025interpretable}
Lv X, Wu Z, Xu J, et al. Interpretable procedural material graph generation via diffusion models from reference images. \textit{The Visual Computer}, 2025, 41(13): 11195--11205.

\bibitem{zeng2024paint3d}
Zeng X, Chen X, Qi Z, et al. Paint3D: Paint anything 3D with lighting-immersive texture generation. \textit{Proceedings of the IEEE/CVF Conference on Computer Vision and Pattern Recognition}, 2024.

\bibitem{cheng2024mvpaint}
Cheng W, Mu J, Zeng X, et al. MVPaint: Synchronized multi-view diffusion for painting anything 3D. \textit{Proceedings of the IEEE/CVF Conference on Computer Vision and Pattern Recognition}, 2025.

\bibitem{lai2025hunyuan3d}
Zhao Z, Lai Z, Lin Q, et al. Hunyuan3D 2.0: Scaling diffusion models for high resolution textured 3D assets generation. arXiv:2501.12202, 2025.

\bibitem{huang2024mvadapter}
Huang Z, et al. MV-Adapter: Multi-view consistent image generation made easy. \textit{Proceedings of the IEEE/CVF International Conference on Computer Vision}, 2025.

\bibitem{tang2023mvdiffusion}
Tang S, Zhang F, Chen J, Wang P, Furukawa Y. MVDiffusion: Enabling holistic multi-view image generation with correspondence-aware diffusion. \textit{Advances in Neural Information Processing Systems}, 2023.

\bibitem{liu2024synchronized}
Liu Y X, Xie M S, Liu H Y, Wong T T. Text-guided texturing by synchronized multi-view diffusion. \textit{SIGGRAPH Asia 2024 Conference Papers}, 2024.

\bibitem{xu2026wondertex}
Xu Q, Han X G, Zhang L. WonderTex: Consistent-and-seamless texture generation with text-guided multi-view image diffusion models. \textit{IEEE Transactions on Visualization and Computer Graphics}, 2026, 32(7): 5815--5825.

\bibitem{mou2023t2iadapter}
Mou C, Wang X, Xie L, et al. T2I-Adapter: Learning adapters to dig out more controllable ability for text-to-image diffusion models. \textit{Proceedings of the AAAI Conference on Artificial Intelligence}, 2024.

\bibitem{zhang2023controlnet}
Zhang L, Rao A, Agrawala M. Adding conditional control to text-to-image diffusion models. \textit{Proceedings of the IEEE/CVF International Conference on Computer Vision}, 2023.

\bibitem{ye2023ipadapter}
Ye H, Zhang J, Liu S, Han X, Yang W. IP-Adapter: Text compatible image prompt adapter for text-to-image diffusion models. arXiv:2308.06721, 2023.

\bibitem{hu2022lora}
Hu E J, Shen Y, Wallis P, et al. LoRA: Low-rank adaptation of large language models. \textit{International Conference on Learning Representations}, 2022.

\bibitem{liu2024dora}
Liu S Y, Wang C Y, Yin H, et al. DoRA: Weight-decomposed low-rank adaptation. \textit{International Conference on Machine Learning}, 2024.

\bibitem{zhang2023adalora}
Zhang Q, Zuo S, Liang C, et al. AdaLoRA: Adaptive budget allocation for parameter-efficient fine-tuning. \textit{International Conference on Learning Representations}, 2023.

\bibitem{kulkarni2026ssi}
Kulkarni A S. Scheduled Style Injection: Expanding the Style-Content Pareto Frontier in Training-Free Diffusion-Based Style Transfer. \textit{Proceedings of the IEEE/CVF Conference on Computer Vision and Pattern Recognition Workshops (CVPRW), NTIRE Workshop}, 2026.

\bibitem{ho2020ddpm}
Ho J, Jain A, Abbeel P. Denoising diffusion probabilistic models. \textit{Advances in Neural Information Processing Systems}, 2020.

\bibitem{song2021ddim}
Song J, Meng C, Ermon S. Denoising diffusion implicit models. \textit{International Conference on Learning Representations}, 2021.

\bibitem{ho2022cfg}
Ho J, Salimans T. Classifier-free diffusion guidance. \textit{NeurIPS Workshop on Deep Generative Models and Downstream Applications}, 2021. arXiv:2207.12598.

\bibitem{castillo2023adaptive}
Castillo A, Kohler J, P\'erez J C, et al. Adaptive guidance: Training-free acceleration of conditional diffusion models. \textit{Proceedings of the AAAI Conference on Artificial Intelligence}, 2025.

\bibitem{wang2004ssim}
Wang Z, Bovik A C, Sheikh H R, Simoncelli E P. Image quality assessment: From error visibility to structural similarity. \textit{IEEE Transactions on Image Processing}, 2004, 13(4): 600--612.

\bibitem{zhang2018lpips}
Zhang R, Isola P, Efros A A, et al. The unreasonable effectiveness of deep features as a perceptual metric. \textit{CVPR}, 2018.

\bibitem{sad2024apg}
Sadat S, Hilliges O, Weber R M. Eliminating oversaturation and artifacts of high guidance scales in diffusion models. \textit{International Conference on Learning Representations}, 2025.

\end{thebibliography}
\endgroup
\end{document}
